\title{8-bit Optimizers via Block-wise Quantization}

\begin{abstract}
Stateful optimizers maintain gradient statistics over time, e.g., the exponentially smoothed sum (SGD with momentum) or squared sum (Adam) of past gradient values. This state can be used to accelerate optimization compared to plain stochastic gradient descent but uses memory that might otherwise be allocated to model parameters, thereby limiting the maximum size of models trained in practice. In this paper, we develop the first optimizers that use 8-bit statistics while maintaining the performance levels of using 32-bit optimizer states. To overcome the resulting computational, quantization, and stability challenges, we develop block-wise dynamic quantization. Block-wise quantization divides input tensors into smaller blocks that are independently quantized. Each block is processed in parallel across cores, yielding faster optimization and high precision quantization. To maintain stability and performance, we combine block-wise quantization with two additional changes: (1) dynamic quantization, a form of non-linear optimization that is precise for both large and small magnitude values, and (2) a stable embedding layer to reduce gradient variance that comes from the highly non-uniform distribution of input tokens in language models. As a result, our 8-bit optimizers maintain 32-bit performance with a small fraction of the memory footprint on a range of tasks, including 1.5B parameter language modeling, GLUE finetuning, ImageNet classification, WMT'14 machine translation, MoCo v2 contrastive ImageNet pretraining+finetuning, and RoBERTa pretraining, without changes to the original optimizer hyperparameters. We open-source our 8-bit optimizers as a drop-in replacement that only requires a two-line code change.
\end{abstract}

\section{Background}

\subsection{Stateful Optimizers}
\label{section:optim}

An optimizer updates the parameters $\mathbf{w}$ of a neural network by using the gradient of the loss with respect to the weight $\mathbf{g}_t=\frac{\mathbf{\partial L}}{\mathbf{\partial w}}$ at update iteration $t$. Stateful optimizers compute statistics of the gradient with respect to each parameter over time for accelerated optimization. Two of the most commonly used stateful optimizers are Adam \citep{kingma2014adam}, and SGD with momentum \citep{Qian1999OnTM} -- or Momentum for short. Without damping and scaling constants, the update rules of these optimizers are given by:

\begin{equation}
   \text{Momentum}(\mathbf{g}_t,\mathbf{w}_{t-1}, \mathbf{m}_{t-1}) = 
   \begin{cases}
   \mathbf{m}_0 = \mathbf{g}_0 & \text{Initialization} \\
   \mathbf{m}_t = \beta_1\mathbf{m}_{t-1} +\mathbf{g}_t & \text{State 1 update} \\
   \mathbf{w}_t = \mathbf{w}_{t-1} - \alpha\cdot\mathbf{m}_t & \text{Weight update}
   \end{cases}
\end{equation}

\begin{equation}
   \text{Adam}(\mathbf{g}_t,\mathbf{w}_{t-1}, \mathbf{m}_{t-1}, \mathbf{r}_{t-1})= 
     \begin{cases}
   \mathbf{r}_0 = \mathbf{m}_0 = \mathbf{0} & \text{Initialization} \\
   \mathbf{m}_t = \beta_1\mathbf{m}_{t-1} + (1-\beta_1)\mathbf{g}_t & \text{State 1 update} \\
   \mathbf{r}_t = \beta_2\mathbf{r}_{t-1} + (1-\beta_2)\mathbf{g}_t^2 & \text{State 2 update} \\
   \mathbf{w}_t = \mathbf{w}_{t-1} - \alpha\cdot\frac{\mathbf{m}_t}{\sqrt{\mathbf{r}_t}+\epsilon} & \text{Weight update,}
   \end{cases}
\end{equation}

where $\beta_1$ and $\beta_2$ are smoothing constants, $\epsilon$ is a small constant, and $\alpha$ is the learning rate.

For 32-bit states, Momentum and Adam consume 4 and 8 bytes per parameter. That is 4 GB and 8 GB for a 1B parameter model. Our 8-bit non-linear quantization reduces these costs to 1 GB and 2 GB.

\subsection{Non-linear Quantization}
\label{section:nonlinearquant}

Quantization compresses numeric representations to save space at the cost of precision. Quantization is the mapping of a $k$-bit integer to a real element in $D$, that is, ${\mathbf{Q}^{\text{map}}: [0, 2^k-1] \mapsto D}$. For example, the IEEE 32-bit floating point data type maps the indices $0...2^{32}-1$ to the domain \mbox{[-3.4e38,~+3.4e38]}. We use the following notation: ${\mathbf{Q}^{\text{map}}(i) = \mathbf{Q}^{\text{map}}_i = q_i}$, for example  ${\mathbf{Q}^{\text{map}}(2^{31}+131072) = 2.03125}$, for the IEEE 32-bit floating point data type.

To perform general quantization from one data type into another we require three steps. (1) Compute a normalization constant $N$ that transforms the input tensor $\mathbf{T}$ into the range of the domain $D$ of the target quantization data type ${\mathbf{Q}^{\text{map}}}$, (2) for each element of $\mathbf{T}/N$ find the closest corresponding value $q_i$ in the domain $D$, (3) store the index $i$ corresponding to $q_i$ in the quantized output tensor $\mathbf{T}^Q$. To receive the dequantized tensor $\mathbf{T}^D$ we look up the index and denormalize: $\mathbf{T}^D_i = {\mathbf{Q}^{\text{map}}}(\mathbf{T}^Q_i)\cdot N$.

To perform this procedure for dynamic quantization we first normalize into the range [-1, 1] through division by the absolute maximum value: $N = \max(|\mathbf{T}|)$.

Then we find the closest values via a binary search:

\begin{equation}
    \mathbf{T}_i^Q = \argmin_{j=0}^{2^n} |\mathbf{Q}^{\text{map}}_j - \frac{\mathbf{T}_i}{N}|
\end{equation}

\subsection{Dynamic Tree Quantization}
 \label{section:dynamictreequant}

Dynamic Tree quantization \citep{dettmers20168bit} is a method that yields low quantization error for both small and large magnitude values. Unlike data types with fixed exponent and fraction, dynamic tree quantization uses a datatype with a dynamic exponent and fraction that can change with each number. It is made up of four parts, as seen in Figure~\ref{fig:dynamic_exponent}: (1) The first bit of the data type is reserved for a sign. (2) The number of subsequent zero bits indicates the magnitude of the exponent. (3) The first bit that is set to one indicates that all following values are reserved for (4) linear quantization. By moving the indicator bit, numbers can have a large exponent $10^{-7}$ or precision as high as $1/63$. Compared to linear quantization, dynamic tree quantization has better absolute and relative quantization errors for non-uniform distributions. Dynamic tree quantization is strictly defined to quantize numbers in the range [-1.0, 1.0], which is ensured by performing tensor-level absolute max normalization. 

\section{8-bit Optimizers}
\label{section:8bit_optim}

Our 8-bit optimizers have three components: (1) block-wise quantization that isolates outliers and distributes the error more equally over all bits; (2) dynamic quantization, which quantizes both small and large values with high precision; and (3) a stable embedding layer to improve stability during optimization for models with word embeddings.

With these components, performing an optimizer update with 8-bit states is straightforward. We dequantize the 8-bit optimizer states to 32-bit, perform the update, and then quantize the states back to 8-bit for storage. We do this 8-bit to 32-bit conversion element-by-element in registers, which means no slow copies to GPU memory or additional temporary memory are needed to perform quantization and dequantization. For GPUs, this makes 8-bit optimizers faster than regular 32-bit optimizers, as we show in Section~\ref{sec:mainresults}.

\subsection{Block-wise Quantization}

Our block-wise quantization reduces the cost of computing normalization and improves quantization precision by isolating outliers. In order to dynamically quantize a tensor, as defined in Section~\ref{section:nonlinearquant}, we need to normalize the tensor into the range [-1, 1]. Such normalization requires a reduction over the entire tensor, which entails multiple synchronizations across GPU cores. Block-wise dynamic quantization reduces this cost by chunking an input tensor into small blocks of size $B=2048$ and performing normalization independently in each core across this block.

More formally, using the notation introduced in Section~\ref{section:nonlinearquant}, in block-wise quantization, we treat $\mathbf{T}$ as a one-dimensional sequence of elements that we chunk in blocks of size $B$. This means for an input tensor $\mathbf{T}$ with $n$ elements we have $n/B$ blocks. We proceed to compute a normalization constant for each block: $N_b = \max(|\mathbf{T}_b|)$, where $b$ is the index of the block $0..n/B$. With this block-wise normalization constant, each block can be quantized independently:

\begin{equation}
    \mathbf{T}_{bi}^Q = \argmin_{j=0}^{2^n} |\mathbf{Q}^{\text{map}}_j - \frac{\mathbf{T}_{bi}}{N_b}|{\bigg\rvert}_{0<i<B}
\end{equation}

This approach has several advantages, both for stability and efficiency. First, each block normalization can be computed independently. Thus no synchronization between cores is required, and throughput is enhanced.

Secondly, it is also much more robust to outliers in the input tensor. For example, to contrast block-wise and regular quantization, if we create an input tensor with one million elements sampled from the standard normal distribution, we expect less than 1\% of elements of the tensor will be in the range [3, $+\infty)$. However, since we normalize the input tensor into the range \mbox{[-1,1]} this means the maximum values of the distribution determine the range of quantization buckets. This means if the input tensor contains an outlier with magnitude 5, the quantization buckets reserved for numbers between 3 and 5 will mostly go unused since less than 1\% of numbers are in this range. With block-wise quantization, the effect of outliers is limited to a single block. As such, most bits are used effectively in other blocks.

Furthermore, because outliers represent the absolute maximum value in the input tensor, block-wise quantization approximates outlier values without any error. This guarantees that the largest optimizer states, arguably the most important, will always be quantized with full precision. This property makes block-wise dynamic quantization both robust and precise and is essential for good training performance in practice.

\subsection{Dynamic Quantization}

In this work, we extend dynamic tree quantization (Section~\ref{section:dynamictreequant}) for non-signed input tensors by re-purposing the sign bit. Since the second Adam state is strictly positive, the sign bit is not needed. Instead of just removing the sign bit, we opt to extend dynamic tree quantization with a {\it fixed} bit for the fraction. This extension is motivated by the observation that the second Adam state varies around 3-5 orders of magnitude during the training of a language model. In comparison, dynamic tree quantization already has a range of 7 orders of magnitude. We refer to this quantization as {\em dynamic quantization} to distinguish it from dynamic tree quantization in our experiments. A study of additional quantization data types and their performance is detailed in Appendix~\ref{appendix:quant}.

\subsection{Stable Embedding Layer}
\label{section:stable_emb}

Our stable embedding layer is a standard word embedding layer variation \citep{devlin2019bert} designed to ensure stable training for NLP tasks.  This embedding layer supports more aggressive quantization by normalizing the highly non-uniform distribution of inputs to avoid extreme gradient variation. See Appendix~\ref{appendix:embeddings} for a discussion of why commonly adopted embedding layers \citep{ott2019fairseq} are so unstable.

We initialize the Stable Embedding Layer with Xavier uniform initialization \citep{glorot2010understanding} and apply layer normalization \citep{ba2016layer} before adding position embeddings. This method maintains a variance of roughly one both at initialization and during training. Additionally, the uniform distribution initialization has less extreme values than a normal distribution, reducing maximum gradient size. Like \cite{ramesh2021zero}, we find that the stability of training improves significantly if we use 32-bit optimizer states for the embedding layers. This is the only layer that uses 32-bit optimizer states. We still use the standard precision for weights and gradients for the embedding layers -- usually 16-bit. We show in our Ablation Analysis in Section~\ref{section:ablations} that the Stable Embedding Layer is required for stable training. See ablations for the Xavier initialization, layer norm, and 32-bit state components of the Stable Embedding Layer in Appendix~\ref{appendix:stable_ablations}.

\section{8-bit vs 32-bit Optimizer Performance for Common Benchmarks}
\label{sec:mainresults}

\paragraph{Experimental Setup}

We compare the performance of 8-bit optimizers to their 32-bit counterparts on a range of challenging public benchmarks. These benchmarks either use Adam \citep{kingma2014adam}, AdamW \citep{loshchilov2018fixing}, or Momentum \citep{Qian1999OnTM}. 

We do not change any hyperparameters or precision of weights, gradients, and activations/input gradients for each experimental setting compared to the public baseline-- the only change is to replace 32-bit optimizers with 8-bit optimizers. This means that for most experiments, we train in 16-bit mixed-precision \citep{micikevicius2017mixed}. We also compare with Adafactor \citep{shazeer2018adafactor}, with the time-independent formulation for $\beta_2$ \citep{shazeer2018adafactor} -- which is the same formulation used in Adam. We also do not change any hyperparameters for Adafactor.

We report on benchmarks in neural machine translation \citep{ott2018scaling}\footnote{\tiny{\url{https://github.com/pytorch/fairseq/tiny/master/examples/scaling_nmt/README.md}}} trained on WMT'16 \citep{sennrich2016edinburgh} and evaluated on en-de WMT'14 \citep{machavcek2014results}, large-scale language modeling \citep{lewis2021base, brown2020language} and RoBERTa pretraining \citep{liu2019roberta} on English CC-100 + RoBERTa corpus \citep{nagel2016cc,Gokaslan2019OpenWeb,zhu2015aligning,wenzek-etal-2020-ccnet}, finetuning the pretrained masked language model RoBERTa \citep{liu2019roberta}\footnote{\tiny{\url{https://github.com/pytorch/fairseq/blob/master/examples/roberta/README.glue.md}}} on GLUE \citep{wang2018glue}, ResNet-50 v1.5 image classification \citep{he2016deep}\footnote{\tiny{\url{https://github.com/NVIDIA/DeepLearningExamples/tree/master/PyTorch/Classification/ConvNets/}}} on ImageNet-1k \citep{deng2009imagenet}, and Moco v2 contrastive image pretraining and linear finetuning  \citep{chen2020improved}\footnote{\tiny{\url{https://github.com/facebookresearch/moco}}} on ImageNet-1k \citep{deng2009imagenet}. 

We use the stable embedding layer for all NLP tasks except for finetuning on GLUE. Beyond this, we follow the exact experimental setup outlined in the referenced papers and codebases. We consistently report replication results for each benchmark with public codebases and report median accuracy, perplexity, or BLEU over ten random seeds for GLUE, three random seeds for others tasks, and a single random seed for large scale language modeling. While it is standard to report means and standard errors on some tasks, others use median performance. We opted to report medians for all tasks for consistency.

\paragraph{Results}

In Table~\ref{tbl:mainresults}, we see that 8-bit optimizers match replicated 32-bit performance for all tasks. While Adafactor is competitive with 8-bit Adam, 8-bit Adam uses less memory and provides faster optimization. Our 8-bit optimizers save up to 8.5 GB of GPU memory for our largest 1.5B parameter language model and 2.0 GB for RoBERTa. Thus, 8-bit optimizers maintain performance and improve accessibility to the finetuning of large models for those that cannot afford GPUs with large memory buffers. We show models that are now accessible with smaller GPUs in Table~\ref{tbl:hfmodels}. A breakdown of individual dataset results on GLUE can be found in Appendix~\ref{appendix:glue}).

The broad range of tasks and competitive results demonstrate that 8-bit optimizers are a robust and effective replacement for 32-bit optimizers, do not require any additional changes in hyperparameters, and save a significant amount of memory while speeding up training slightly.

\section{Analysis}
\label{sec:ablations}

We analyze our method in two ways. First, we ablate all 8-bit optimizer components and show that they are necessary for good performance. Second, we look at the sensitivity to hyperparameters compared to 32-bit Adam and show that 8-bit Adam with block-wise dynamic quantization is a reliable replacement that does not require further hyperparameter tuning.

\paragraph{Experimental Setup} We perform our analysis on a strong 32-bit Adam baseline for language modeling with transformers \citep{vaswani2017attention}. We subsample from the RoBERTa corpus \citep{liu2019roberta} which consists of the English sub-datasets: Books \citep{zhu2015aligning}, Stories \citep{trinh2018simple}, OpenWebText-1 \citep{Gokaslan2019OpenWeb}, Wikipedia, and CC-News \citep{nagel2016cc}. We use a 50k token BPE encoded vocabulary \citep{sennrich2015neural}. We find the best 2-GPU-day transformer baseline for 32-bit Adam with multiple hyperparameter searches that take in a total of 440 GPU days. Key hyperparameters include 10 layers with a model dimension of 1024, a fully connected hidden dimension of 8192, 16 heads, and input sub-sequences with a length of 512 tokens each. The final model has 209m parameters.

\paragraph{Ablation Analysis}
\label{section:ablations}

For the ablation analysis, we compare small and large-scale language modeling perplexity and training stability against a 32-bit Adam baseline. We ablate components individually and include combinations of methods that highlight their interactions. The baseline method uses linear quantization, and we add dynamic quantization, block-wise quantization, and the stable embedding layer to demonstrate their effect. To test optimization stability for small-scale language modeling, we run each setting with different hyperparameters and report median performance across all successful runs. A successful run is a run that does not crash due to exploding gradients or diverges in the loss. We use the hyperparameters $\epsilon$ \{1e-8, 1e-7, 1e-6\}, $\beta_1$ \{0.90, 0.87, 0.93\}, $\beta_2$ \{0.999, 0.99, 0.98\} and small changes in learning rates. We also include some partial ablations for large-scale models beyond 1B parameters. In the large-scale setting, we run several seeds with the same hyperparameters. We use a single seed for 32-bit Adam, five seeds for 8-bit Adam at 1.3B parameters, and a single seed for 8-bit Adam at 1.5B parameters.\footnote{We chose not to do the full ablations with such large models because each training run takes one GPU year.} Results are shown in Table~\ref{tbl:ablations}. 

The Ablations show that dynamic quantization, block-wise quantization, and the stable embedding layer are critical for either performance or stability. In addition, block-wise quantization is critical for large-scale language model stability.

\paragraph{Sensitivity Analysis}

We compare the perplexity of 32-bit Adam vs 8-bit Adam + Stable Embedding as we change the optimizer hyperparameters: learning rate, betas, and $\epsilon$. We change each hyperparameter individually from the baseline hyperparameters $\beta_1$=0.9, $\beta_2$=0.995, $\epsilon$=1e-7, and lr=0.0163 and run two random seeds for both 8-bit and 32-bit Adam for each setting. If 8-bit Adam is perfectly insensitive to hyperparameters compared to 32-bit Adam, we would expect the same constant offset in performance for any hyperparameter combination. The results can be seen in Figure~\ref{fig:sensitivity}. The results show a relatively steady gap between 8-bit and 32-bit Adam, suggesting that 8-bit Adam does not require any further hyperparameter tuning compared to 32-bit Adam.

\section{Related Work}

\paragraph{Compressing \& Distributing Optimizer States} While 16-bit Adam has been used in several publications, the stability of 16-bit Adam was first explicitly studied for a text-to-image generation model DALL-E \citep{ramesh2021zero}. They show that a stable embedding layer, tensor-wise scaling constants for both Adam states, and multiple loss scaling blocks are critical to achieving stability during training. Our work reduces the memory footprint of Adam further, from 16 to 8-bit. In addition, we achieve stability by developing new training procedures and non-linear quantization, both of which complement previous developments.

Adafactor \citep{shazeer2018adafactor} uses a different strategy to save memory. All optimizer states are still 32-bit, but the second Adam state is factorized by a row-column outer product resulting in a comparable memory footprint to 16-bit Adam. Alternatively, Adafactor can also be used without using the first moment ($\beta_1 = 0.0$) \citep{shazeer2018adafactor}. This version is as memory efficient as 8-bit Adam, but unlike 8-bit Adam, hyperparameters for this Adafactor variant need to be re-tuned to achieve good performance. We compare 8-bit Adam with Adafactor $\beta_1 > 0.0$ in our experiments. 

AdaGrad \citep{duchi2011adaptive} adapts the gradient with aggregate training statistics over the entire training run. AdaGrad that uses only the main diagonal as optimizer state and extensions of AdaGrad such as SM3 \citep{anil2019sm3} and extreme tensoring \citep{chen2020extremetensor} can be more efficient than 8-bit Adam. We include some initial comparison with AdaGrad in Appendix~\ref{appendix:adagrad}.

Optimizer sharding \citep{rajbhandari2020zero} splits optimizer states across multiple accelerators such as GPUs/TPUs. While very effective, it can only be used if multiple accelerators are available and data parallelism is used. Optimizer sharding can also have significant communication overhead \citep{rajbhandari2021zero}. Our 8-bit optimizers work with all kinds of parallelism. They can also complement optimizer sharding, as they reduce communication overhead by 75\%. 

\paragraph{General Memory Reduction Techniques} Other complementary methods for efficient training can be either distributed or local. Distributed approaches spread out the memory of a model across several accelerators such as GPUs/TPUs. Such approaches are model parallelism \citep{krizhevsky2009learning}, pipeline parallelism \citep{krizhevsky2009learning,huang2018gpipe,harlap2018pipedream}, and operator parallelism \citep{lepikhin2020gshard}. These approaches are useful if one has multiple accelerators available. Our 8-bit optimizers are useful for both single and multiple devices.

Local approaches work for a single accelerator. They include gradient checkpointing \citep{chen2016training}, reversible residual connections \citep{gomez2017reversible}, and offloading \citep{pudipeddi2020training,rajbhandari2021zero}. All these methods save memory at the cost of increased computational or communication costs. Our 8-bit optimizers reduce the memory footprint of the model while maintaining 32-bit training speed.

\paragraph{Quantization Methods and Data Types}

While our work is the first to apply 8-bit quantization to optimizer statistics, quantization for neural network model compression, training, and inference are well-studied problems. One of the most common formats of 8-bit quantization is to use data types composed of static sign, exponent, and fraction bits. The most common combination is 5 bits for the exponent and 2 bits for the fraction \citep{wang2018training8bit,sun2019hybrid8bit,cambier2020shiftsqueeze,mellempudi2019bit8} with either no normalization or min-max normalization. These data types offer high precision for small magnitude values but have large errors for large magnitude values since only 2 bits are assigned to the fraction. Other methods improve quantization through soft constraints \citep{li2021endtoend} or more general uniform affine quantizations \citep{brevitas}.

Data types lower than 8-bit are usually used to prepare a model for deployment, and the main focus is on improving network inference speed and memory footprint rather than maintaining accuracy. There are methods that use 1-bit \citep{courbariaux2016bit1,Rastegari2016xnor,courbariaux2015binaryconnect}, 2-bit/3 values \citep{zhu2017ternary,choi2019bit2}, 4-bits \citep{li2019bit4}, more bits \citep{courbariaux2014training}, or a variable amount of bits \citep{gong2019softquant}. See also \cite{qin2020survey1bit} for a survey on binary neural networks. While these low-bit quantization techniques allow for efficient storage, they likely lead to instability when used for optimizer states.

The work most similar to our block-wise quantization is work on Hybrid Block Floating Point (HBFP) \citep{drumond2018hybridblock} which uses a 24-bit fraction data type with a separate exponent for each tile in matrix multiplication to perform 24-bit matrix multiplication. However, unlike HBFP, block-wise dynamic quantization has the advantage of having both block-wise normalization {\it and} a dynamic exponent for each number. This allows for a much broader range of important values since optimizer state values vary by about 5 orders of magnitude. Furthermore, unlike HBFP, block-wise quantization approximates the maximum magnitude values within each block without any quantization error, which is critical for optimization stability, particularly for large networks.

\section{Discussion \& Limitations}

Here we have shown that high precision quantization can yield 8-bit optimizers that maintain 32-bit optimizer performance without requiring any change in hyperparameters. One of the main limitations of our work is that 8-bit optimizers for natural language tasks require a stable embedding layer to be trained to 32-bit performance. On the other hand, we show that 32-bit optimizers also benefit from a stable embedding layer. As such, the stable embedding layer could be seen as a general replacement for other embedding layers.

We show that 8-bit optimizers reduce the memory footprint and accelerate optimization on a wide range of tasks. However, since 8-bit optimizers reduce only the memory footprint proportional to the number of parameters, models that use large amounts of activation memory, such as convolutional networks, have few benefits from using 8-bit optimizers. Thus, 8-bit optimizers are most beneficial for training or finetuning models with many parameters on highly memory-constrained GPUs.

Furthermore, there remain sources of instability that, to our knowledge, are not well understood. For example, we observed that models with over 1B parameters often have hard systemic divergence, where many parameters simultaneously cause exploding gradients. In other cases, a single parameter among those 1B parameters assumed a value too large, caused an exploding gradient, and led to a cascade of instability. It might be that this rare cascading instability is related to the phenomena where instability disappears after reloading a model checkpoint and rolling a new random seed -- a method standard for training huge models. Cascading instability might also be related to the observation that the larger a model is, the more unstable it becomes. For 8-bit optimizers, handling outliers through block-wise quantization and the stable embedding layer was key for stability. We hypothesize that that extreme outliers are related to cascading instability. If such phenomena were better understood, it could lead to better 8-bit optimizers and stable training in general.

\end{document}
